\lecture{4.1}{From Identification to Estimation}
\label{lec:4-1}

\subsection{Observed Data Structure: $O=(W,A,Y)$}

In the previous chapter, the average treatment effect was identified under consistency,
conditional randomization, and positivity. Identification gave formulas involving only
the observed-data distribution. Estimation begins after that step: we now have a
sample and must turn the identified statistical parameter into a numerical estimate.

Throughout this chapter the observed data are
\[
O_i=(W_i,A_i,Y_i), \qquad i=1,\ldots,n,
\]
where the observations are assumed to be independent and identically distributed:
\[
O_1,\ldots,O_n \iid P.
\]
Here $W$ is a vector of pre-treatment covariates, $A\in\{0,1\}$ is treatment, and
$Y$ is the outcome. The phrase ``pre-treatment'' is important. Variables affected by
treatment are not ordinary confounders for the effect of $A$ on $Y$ in a point-treatment
analysis, because adjusting for them can block part of the causal effect or create
collider bias.

\begin{assumptionbox}
The estimators in this chapter are estimators of causal effects only when the
identification assumptions from Chapter 3 are scientifically credible:
\[
Y(a)=Y \text{ if } A=a,\qquad
\Ya{a}\ci A\given W,\qquad
\Prob(A=a\given W=w)>0
\]
for all treatment levels $a$ and all covariate values $w$ in the target population.
Without these assumptions, the same formulas estimate well-defined statistical
parameters, but those parameters need not equal counterfactual means.
\end{assumptionbox}

\subsection{Outcome Regression $\Qbar(a,w)$}

\begin{definition}[Outcome regression]
The outcome regression is the conditional mean of the observed outcome given
treatment and covariates:
\[
\Qbar(a,w)=\E(Y\given A=a,W=w).
\]
\end{definition}

This function answers an observed-data question: among individuals who actually
received treatment level $a$ and had covariates $W=w$, what was the mean outcome?
Under the identification assumptions, this observed subgroup mean can be used as a
building block for the mean counterfactual outcome under treatment level $a$.

For a binary outcome, $\Qbar(a,w)$ is a conditional risk:
\[
\Qbar(a,w)=\Prob(Y=1\given A=a,W=w).
\]
For a continuous outcome, it is a conditional average. For an ordinal or count outcome,
it is still a conditional mean, although the mean may not fully summarize the whole
conditional distribution.

\subsection{Propensity Score $g_a(w)$}

\begin{definition}[Propensity score]
For a binary treatment, the treatment-specific propensity score is
\[
g_a(w)=\Prob(A=a\given W=w).
\]
The usual propensity score is $g_1(w)=\Prob(A=1\given W=w)$, and
$g_0(w)=1-g_1(w)$.
\end{definition}

The propensity score describes how treatment is assigned in the observed population.
It is not a causal effect. Instead, it measures the treatment probability among people
with covariate value $w$. In inverse probability weighting, individuals who received a
treatment level that was unlikely for their covariates receive larger weights. These
weights compensate for the unequal representation of covariate strata across treatment
groups.

\subsection{Plug-In Estimation}

The g-computation identification formula for the counterfactual mean is
\[
\psi(a)=\E\{\Ya{a}\}=\E\{\Qbar(a,W)\}.
\]
This formula contains two features of the observed-data distribution:
the conditional mean function $\Qbar(a,w)$ and the marginal distribution of $W$.
A plug-in estimator replaces each unknown population quantity by an estimate from
the data.

\begin{definition}[Plug-in estimator]
A plug-in estimator is obtained by writing the target parameter as a function of the
observed-data distribution and then substituting empirical or model-based estimates
for the unknown components of that distribution.
\end{definition}

If $\Qbar_n(a,w)$ estimates $\Qbar(a,w)$, the natural g-computation plug-in estimator is
\[
\psi_n^{\mathrm{gcomp}}(a)
=
\frac{1}{n}\sum_{i=1}^n \Qbar_n(a,W_i).
\]
This expression averages predicted outcomes under treatment level $a$ over the
observed covariate distribution. The observed covariates $W_1,\ldots,W_n$ are used
as an empirical approximation to the target population distribution of $W$.

\begin{figure}[H]
\centering
\begin{tikzpicture}[
  scale=.88,
  transform shape,
  node distance=1.25cm,
  box/.style={draw, rounded corners=2pt, minimum width=2.8cm, minimum height=.85cm},
  arr/.style={-{Latex[length=2.5mm]}, thick}
]
\node[box] (data) {Observed data\\ $O_i=(W_i,A_i,Y_i)$};
\node[box, right=of data] (fit) {Fit nuisance\\ functions};
\node[box, right=of fit] (plugin) {Plug into\\ identified formula};
\node[box, right=of plugin] (ate) {Estimate\\ $\psi(1)-\psi(0)$};
\draw[arr] (data) -- (fit);
\draw[arr] (fit) -- node[above] {$\Qbar_n,\ g_n$} (plugin);
\draw[arr] (plugin) -- (ate);
\end{tikzpicture}
\caption{Estimation turns the identified observed-data parameter into a sample quantity.}
\label{fig:ch4-estimation-flow}
\end{figure}

\subsection{Estimating Counterfactual Means}

There are two basic plug-in strategies emphasized in this chapter.

The first uses the outcome regression:
\[
\psi_n^{\mathrm{gcomp}}(a)=\frac{1}{n}\sum_{i=1}^n \Qbar_n(a,W_i).
\]
For each observed individual, we keep $W_i$ fixed and replace the observed treatment
by the hypothetical treatment value $a$ inside the fitted outcome regression. The
resulting predicted values are then averaged.

The second uses inverse probability weighting. If $g_{n,a}(w)$ estimates
$g_a(w)=\Prob(A=a\given W=w)$, then
\[
\psi_n^{\IPTW}(a)
=
\frac{1}{n}\sum_{i=1}^n
\frac{\ind(A_i=a)}{g_{n,a}(W_i)}Y_i.
\]
Only people who actually received treatment level $a$ contribute to the estimate of
$\E\{\Ya{a}\}$, but their outcomes are weighted by the inverse of their estimated
probability of receiving that treatment.

\begin{keyidea}
G-computation estimates counterfactual means by predicting everyone's outcome under
each treatment level. IPTW estimates counterfactual means by reweighting the people
who actually received each treatment level so that, after weighting, they represent the
full target population.
\end{keyidea}

\subsection{Estimating the ATE}

The average treatment effect is
\[
\ATE=\E\{\Yone\}-\E\{\Yzero\}=\psi(1)-\psi(0).
\]
The corresponding g-computation estimator is
\[
\widehat{\ATE}_{\mathrm{gcomp}}
=
\frac{1}{n}\sum_{i=1}^n
\left\{
\Qbar_n(1,W_i)-\Qbar_n(0,W_i)
\right\}.
\]
The corresponding IPTW estimator is
\[
\widehat{\ATE}_{\IPTW}
=
\frac{1}{n}\sum_{i=1}^n
\left\{
\frac{\ind(A_i=1)}{g_{n,1}(W_i)}Y_i
-
\frac{\ind(A_i=0)}{g_{n,0}(W_i)}Y_i
\right\}.
\]

\begin{estimationbox}
Basic estimation workflow for the ATE:
\begin{enumerate}
\item Choose whether to estimate the counterfactual means using the outcome
regression, the propensity score, or both in a later doubly robust method.
\item Estimate the nuisance function required by the chosen method.
\item Construct $\psi_n(1)$ and $\psi_n(0)$ using the identified formula.
\item Report $\widehat{\ATE}=\psi_n(1)-\psi_n(0)$.
\item Quantify uncertainty using an appropriate variance estimator or resampling
method such as the bootstrap.
\end{enumerate}
\end{estimationbox}
